\thispagestyle{empty}
\begin{center}
\vspace*{1.1in}
{\sffamily\fontsize{40}{44}\selectfont\bfseries\color{stblue} Smalltalk-2026}

\vspace{10pt}
{\color{strule}\rule{0.62\textwidth}{1pt}}

\vspace{14pt}
{\sffamily\Large\color{stink} A Tutorial}

\vspace{6pt}
{\sffamily\normalsize\color{stgrey} the language, the library, the image, and the cores}

\vfill

{\small\color{stgrey}
\url{https://github.com/blakemcbride/Smalltalk-2026}
}

\vspace{18pt}
\end{center}

\clearpage
\thispagestyle{empty}
\vspace*{\fill}
\noindent
\begin{minipage}{\textwidth}
\small\color{stgrey}
\noindent
Copyright \copyright{} 2026 Blake McBride. This document is distributed under
the BSD 2-Clause licence, as is the rest of \ct{doc/}.

\vspace{8pt}\noindent
Smalltalk-2026 is not all one licence, and the distinction matters. The 1983
class library in \ct{sources/} is MIT, from \emph{markbush/Smalltalk-80-Sources},
and is vendored unedited. The packages under \ct{pharo/} are MIT with parts
under Apache-2.0, each carrying its own notice and a \ct{PROVENANCE.md}. The
rasterised face in \ct{src/gfx/font_face.c} derives from Inter, under the SIL
Open Font License 1.1. See \ct{doc/LICENSING.md}.

\vspace{8pt}\noindent
Every Smalltalk expression in this book whose answer is shown was run against
a live image built from this source tree. The examples are checked in as
\ct{tutorial/examples/TutorialExamples.class.st} and \ct{make verify} runs
them. If the book and the system ever disagree, that command is the referee.
\end{minipage}
\vspace*{\fill}
